\section{Requisitos do Sistema e Diretrizes Operacionais}

A especificação de requisitos em sistemas hospitalares de missão crítica demanda rigor metodológico, critérios objetivos de aceitação e rastreabilidade integral. Falhas funcionais ou ambiguidades conceituais em saúde comprometem diretamente a integridade física de pacientes, além de acarretarem sanções regulatórias e perdas financeiras às instituições \cite{sommerville:2019, pressman:2016}.

O Vitalis v2.0 combina regras de negócio determinísticas com serviços de inteligência artificial clínica. Para disciplinar a coexistência desses componentes, os requisitos adotam a priorização MoSCoW:
\begin{itemize}[leftmargin=*]
    \item \textbf{Essencial (\textit{Must have})}: Funcionalidades indispensáveis para a continuidade assistencial, segurança do paciente, conformidade com o Conselho Federal de Medicina (CFM) \cite{cfm_prontuario:2007, cfm_digital:2021} e preservação da integridade transacional.
    \item \textbf{Importante (\textit{Should have})}: Capacidades de alto valor que elevam a produtividade médica, previnem glosas e otimizam a ocupação de leitos, contornáveis temporariamente por procedimentos operacionais padrão.
    \item \textbf{Desejável (\textit{Could have})}: Recursos de conveniência analítica, agendamentos avançados e exportações complementares, implementados conforme a disponibilidade do cronograma.
\end{itemize}

Esta seção estrutura os 22 Requisitos Funcionais (RF), as 12 Regras de Negócio (RN), os 7 Requisitos Não Funcionais (RNF) e a Matriz de Conformidade Regulatória do Vitalis v2.0.

% ==============================================================================
% 5.1 REQUISITOS FUNCIONAIS
% ==============================================================================
\subsection{Requisitos Funcionais (RF)}

Os Requisitos Funcionais abrangem desde o acolhimento do paciente na recepção até o faturamento de encerramento da conta, integrando processamento de linguagem natural e modelos preditivos à rotina clínica. A Tabela \ref{tab:rf_vitalis} consolida o conjunto de requisitos funcionais com seus respectivos códigos, descrições técnicas, critérios mensuráveis de aceitação e prioridades.

\begingroup
\small
\begin{xltabular}{\textwidth}{l >{\raggedright\arraybackslash}p{4.4cm} >{\raggedright\arraybackslash}X c}
\caption{Matriz Detalhada de Requisitos Funcionais do Vitalis v2.0} \label{tab:rf_vitalis} \\
\toprule
\textbf{Código} & \textbf{Descrição do Requisito} & \textbf{Critério de Aceitação Técnico} & \textbf{Prioridade} \\
\midrule
\endfirsthead

\multicolumn{4}{c}%
{{\tablename\ \thetable{} -- Continuação da página anterior}} \\
\toprule
\textbf{Código} & \textbf{Descrição do Requisito} & \textbf{Critério de Aceitação Técnico} & \textbf{Prioridade} \\
\midrule
\endhead

\midrule
\multicolumn{4}{r}{{Continua na próxima página \ldots}} \\
\bottomrule
\endfoot

\bottomrule
\endlastfoot

\multicolumn{4}{l}{\textbf{\color{vitalisblue}Bloco 1: Acolhimento, Triagem e Admissão}} \\
\midrule
\textbf{RF-001} &
Cadastro de pacientes com integração ao barramento nacional CNS/CPF e validação cadastral. &
Validação de duplicidade cadastral e conferência de dígitos verificadores em tempo $< 1{,}0\,\text{s}$. &
Essencial \\
\addlinespace
\textbf{RF-002} &
Cálculo de prioridade clínica por algoritmo de Manchester com validação hemodinâmica determinística. &
Atribuição determinística da cor de risco (Vermelho a Azul) conforme sinais vitais registrados sem intervenção discricionária do software. &
Essencial \\
\addlinespace
\textbf{RF-003} &
Transcrição por voz e extração de sintomas clínicos e queixas principais via API de IA. &
Taxa de erro de palavras ($\text{WER} < 5\%$) para vocabulário médico em língua portuguesa e latência de processamento $< 2{,}0\,\text{s}$. &
Importante \\
\midrule

\multicolumn{4}{l}{\textbf{\color{vitalisblue}Bloco 2: Monitoramento Clínico, Leitos e Sepse}} \\
\midrule
\textbf{RF-004} &
Cálculo contínuo de deterioração hemodinâmica e risco de sepse (escores qSOFA e NEWS2) por IA clínica. &
Atualização automática do escore em tela a cada novo sinal vital lançado, com notificação de alerta crítico em $< 1{,}0\,\text{s}$. &
Essencial \\
\addlinespace
\textbf{RF-005} &
Gestão em tempo real de status de leitos hospitalares (Livre, Ocupado, Limpeza, Manutenção). &
Propagação reativa do estado de cada leito para todos os terminais conectados via WebSockets em tempo $< 200\,\text{ms}$. &
Essencial \\
\addlinespace
\textbf{RF-006} &
Alocação preditiva de leitos baseada em estimativa de tempo de permanência (\textit{Length of Stay}, LOS). &
Geração de recomendação de ala e tipo de leito compatível com o perfil clínico em $< 1{,}5\,\text{s}$ no momento da indicação de internação. &
Importante \\
\midrule

\multicolumn{4}{l}{\textbf{\color{vitalisblue}Bloco 3: Prescrição Médica e Segurança Farmacológica}} \\
\midrule
\textbf{RF-007} &
Registro de prescrição médica digital com checagem cruzada automática de alergias e comorbidades. &
Bloqueio imediato da seleção de fármacos contraindicados no cadastro de alergias do paciente, com registro do veto no prontuário. &
Essencial \\
\addlinespace
\textbf{RF-008} &
Copiloto clínico de IA para validação posológica e detecção de interações medicamentosas via RAG. &
Exibição de análise contextual fundamentada em diretrizes farmacêuticas oficiais em tempo $< 1{,}5\,\text{s}$ durante a digitação. &
Importante \\
\addlinespace
\textbf{RF-009} &
Assinatura digital padrão ICP-Brasil em prontuários, prescrições e laudos (PAdES / PKCS\#7). &
Geração de artefato assinado com carimbo do tempo criptográfico, verificável de forma autônoma fora da plataforma. &
Essencial \\
\addlinespace
\textbf{RF-010} &
Módulo de aprazamento de enfermagem e dupla checagem biométrica para fármacos de alta vigilância. &
Confirmação obrigatória por dois profissionais de enfermagem distintos via leitura biométrica ou senha antes da liberação do lote no leito. &
Essencial \\
\midrule

\multicolumn{4}{l}{\textbf{\color{vitalisblue}Bloco 4: Documentação Clínica, Faturamento e TISS/TUSS}} \\
\midrule
\textbf{RF-011} &
Geração automatizada de minutas de sumários de alta e evolução clínica estruturada por IA. &
Elaboração do rascunho textual estruturado em $< 3{,}0\,\text{s}$, mantendo a obrigatoriedade de revisão e assinatura médica. &
Importante \\
\addlinespace
\textbf{RF-012} &
Emissão e exportação automatizada de guias de faturamento hospitalar no padrão TISS/TUSS da ANS. &
Geração de lotes XML compatíveis com os esquemas XSD oficiais da ANS, com validação sintática pré-envio sem erros de formatação. &
Essencial \\
\addlinespace
\textbf{RF-013} &
Auditoria preventiva de faturamento assistida por IA para identificação de inconsistências e prevenção de glosas. &
Confronto entre registros de procedimentos, medicamentos administrados e regras contratuais com marcação de pendências antes da guia XML. &
Importante \\
\midrule

\multicolumn{4}{l}{\textbf{\color{vitalisblue}Bloco 5: Operação Hospitalar e Farmácia Satélite}} \\
\midrule
\textbf{RF-014} &
Painel gerencial de indicadores em tempo real para monitoramento de tempo de espera e giro de leitos. &
Atualização assíncrona dos cartões métricos do dashboard a cada $30\,\text{s}$ sem recarga da página web. &
Importante \\
\addlinespace
\textbf{RF-015} &
Agendamento inteligente de exames complementares e alocação de salas diagnósticas. &
Identificação dinâmica de conflitos de horário, disponibilidade de equipamentos e tempo de preparo prévio do paciente. &
Desejável \\
\addlinespace
\textbf{RF-016} &
Gestão de estoque de farmácias satélites com baixa automática sincronizada à checagem de beira de leito. &
Decremento imediato no saldo de estoque e disparo automático de ordem de reposição ao atingir o ponto de ressuprimento. &
Essencial \\
\addlinespace
\textbf{RF-017} &
Protocolo de encerramento definitivo de prontuário, alta hospitalar e registro de óbito. &
Bloqueio imediato de escrituração retroativa após validação médica e geração do lote contábil de fechamento do atendimento. &
Essencial \\
\midrule

\multicolumn{4}{l}{\textbf{\color{vitalisblue}Bloco 6: Governança de IA, Segurança e Resiliência}} \\
\midrule
\textbf{RF-018} &
Pipeline de sanitização e mascaramento de dados identificáveis (PII) antes do envio a provedores externos de IA. &
Remoção integral de CPF, CNS, nome completo, telefone e dados de contato de qualquer carga útil transmitida para fora do perímetro seguro. &
Essencial \\
\addlinespace
\textbf{RF-019} &
Rastreabilidade, auditoria e log imutável de todas as inferências, prompts e respostas de modelos de IA. &
Persistência estruturada em trilha append-only do modelo utilizado, temperatura, tokens, versão e feedback de aceite/recusa médica. &
Essencial \\
\addlinespace
\textbf{RF-020} &
Mecanismo de chaveamento de contingência (\textit{Circuit Breaker}) e fallback determinístico para serviços de IA. &
Comutação transparente para rotinas lógicas tradicionais em caso de latência de IA $> 2{,}0\,\text{s}$ ou erro HTTP $5\text{xx}$, sem travar o sistema. &
Essencial \\
\addlinespace
\textbf{RF-021} &
Sistema de emissão de alertas multicanal (visual em tela, sinal sonoro e push) para intercorrências críticas. &
Disparo sonoro imediato e trava visual vermelha no painel central de enfermagem quando um paciente for classificado como Vermelho. &
Essencial \\
\addlinespace
\textbf{RF-022} &
Exportação estruturada de relatórios epidemiológicos e administrativos em formatos PDF, CSV e JSON HL7 FHIR. &
Disponibilização de relatórios filtrados por período, especialidade e desfecho clínico gerados em tempo $< 5{,}0\,\text{s}$. &
Desejável \\
\end{xltabular}
\endgroup

% ==============================================================================
% 5.2 REGRAS DE NEGÓCIO
% ==============================================================================
\subsection{Regras de Negócio e Diretrizes Operacionais (RN)}

As regras de negócio delimitam o funcionamento do Vitalis v2.0, convertendo normativas clínicas, resoluções de conselhos profissionais e imperativos legais em restrições lógicas e transacionais invioláveis.

\begin{enumerate}[label=\textbf{RN-\arabic*:}, leftmargin=*]
    \item \textbf{Prioridade Absoluta para a Classificação Vermelha (Emergência)}:
    \begin{itemize}[noitemsep]
        \item \textit{Declaração}: Todo paciente classificado com a cor Vermelha no Protocolo de Manchester possui prioridade assistencial absoluta sobre qualquer outro atendimento na fila médica.
        \item \textit{Ação Sistêmica}: O sistema realoca instantaneamente o paciente para o topo da lista de espera do pronto atendimento, desconsidera a ordem cronológica de chegada, emite sinal sonoro de alerta no posto de enfermagem e trava a chamada de casos de menor gravidade até que a equipe assuma o atendimento crítico.
        \item \textit{Fundamentação}: Diretrizes do Protocolo Internacional de Manchester \cite{manchester:2014}, que determinam tempo zero de espera para risco iminente de morte.
    \end{itemize}

    \item \textbf{Princípio do Human-in-the-Loop em Decisões de IA}:
    \begin{itemize}[noitemsep]
        \item \textit{Declaração}: Nenhuma sugestão diagnóstica, escore de deterioração, cálculo posológico ou minuta textual gerada por IA pode ser executada ou registrada como ato médico definitivo de forma autônoma.
        \item \textit{Ação Sistêmica}: O sistema exige revisão ativa, confirmação expressa, edição ou recusa por profissional de saúde com registro ativo em conselho (CRM ou COREN), antes que qualquer saída de IA produza efeito assistencial.
        \item \textit{Fundamentação}: Resolução CFM nº 2.299/2021 \cite{cfm_digital:2021} e princípios de responsabilidade clínica em inteligência artificial médica \cite{topol:2019}.
    \end{itemize}

    \item \textbf{Desidentificação e Anonimização de Dados Sensíveis}:
    \begin{itemize}[noitemsep]
        \item \textit{Declaração}: Dados clínicos transmitidos a serviços de inferência de IA não podem conter elementos que permitam a reidentificação direta do paciente no ambiente externo.
        \item \textit{Ação Sistêmica}: O pipeline de pré-processamento executa tokenização criptográfica e substituição por entidades sintéticas (substituindo nomes próprios por pseudônimos e suprimindo números de CPF, RG, CNS, telefones e endereços) antes de qualquer envio de carga útil via rede.
        \item \textit{Fundamentação}: Artigos 5º, 11 e 13 da Lei Geral de Proteção de Dados Pessoais (LGPD, Lei nº 13.709/2018) \cite{brasil_lgpd:2018}.
    \end{itemize}

    \item \textbf{Resiliência Operacional e Chaveamento Circuit Breaker}:
    \begin{itemize}[noitemsep]
        \item \textit{Declaração}: A indisponibilidade, lentidão ou falha de componentes de inteligência artificial não pode paralisar as rotinas assistenciais primárias do hospital.
        \item \textit{Ação Sistêmica}: O sistema monitora o tempo de resposta das chamadas de IA. Caso o tempo exceda o limite de $2{,}0\,\text{s}$ ou a taxa de erros HTTP $5\text{xx}$ ultrapasse $5\%$ em janela móvel de 1 minuto, o padrão \textit{Circuit Breaker} abre o circuito, comutando para algoritmos determinísticos e mantendo a interface disponível para preenchimento manual imediato.
        \item \textit{Fundamentação}: Requisitos de engenharia de software para sistemas de missão crítica \cite{sommerville:2019}.
    \end{itemize}

    \item \textbf{Vedação de Lançamentos Financeiros Retroativos}:
    \begin{itemize}[noitemsep]
        \item \textit{Declaração}: Encerrado e faturado contabilmente um episódio de atendimento hospitalar, é terminantemente proibida a inclusão, alteração ou exclusão de itens de consumo, procedimentos médicos ou taxas de diária.
        \item \textit{Ação Sistêmica}: O banco de dados bloqueia alterações na conta através de restrições transacionais e gatilhos de imutabilidade. Qualquer retificação necessária após o fechamento contábil deve ser realizada por meio de guia suplementar de ajuste com justificativa da auditoria.
        \item \textit{Fundamentação}: Padrão TISS da Agência Nacional de Saúde Suplementar (ANS) \cite{ans_tiss:2023}.
    \end{itemize}

    \item \textbf{Sigilo Assistencial, Privilégio Mínimo e Vínculo Ativo}:
    \begin{itemize}[noitemsep]
        \item \textit{Declaração}: O acesso aos dados integrais do Prontuário Eletrônico do Paciente (PEP) é restrito exclusivamente aos profissionais de saúde formalmente vinculados ao episódio de atendimento ativo.
        \item \textit{Ação Sistêmica}: O motor de controle de acesso (RBAC/ABAC) restringe visualizações completas. Tentativas de acesso por profissionais sem vínculo clínico ativo exigem justificativa formal de emergência registrada em trilha de auditoria imutável, notificando o Encarregado de Proteção de Dados (DPO) da instituição.
        \item \textit{Fundamentação}: Código de Ética Médica e Resolução CFM nº 1.821/2007 \cite{cfm_prontuario:2007}.
    \end{itemize}

    \item \textbf{Alocação Exclusiva em Leitos Homologados e Higienizados}:
    \begin{itemize}[noitemsep]
        \item \textit{Declaração}: Nenhum paciente pode ser admitido ou transferido para um leito hospitalar cujo estado não esteja formalmente certificado como ``Livre'' e ``Higienizado''.
        \item \textit{Ação Sistêmica}: A interface de gestão de internação bloqueia a confirmação de transferência caso o leito alvo esteja com status ``Ocupado'', ``Reservado'', ``Em Limpeza'' ou ``Em Manutenção''. A transição para o estado ``Livre'' exige a confirmação digital da equipe de hotelaria hospitalar.
        \item \textit{Fundamentação}: Protocolos de controle de infecção da Comissão de Controle de Infecção Hospitalar (CCIH) e normas sanitárias da ANVISA.
    \end{itemize}

    \item \textbf{Validade Jurídica e Autenticidade de Prescrições}:
    \begin{itemize}[noitemsep]
        \item \textit{Declaração}: Prescrições de medicamentos, ordens cirúrgicas e solicitações de dietas especiais somente adquirem validade assistencial após assinatura digital com certificado no padrão da Infraestrutura de Chaves Públicas Brasileira (ICP-Brasil).
        \item \textit{Ação Sistêmica}: O sistema impede a emissão de ordens de dispensação farmacêutica e o aprazamento de enfermagem enquanto a prescrição constar em status de ``Rascunho'' ou contiver assinatura criptográfica inválida ou expirada.
        \item \textit{Fundamentação}: Lei Federal nº 14.063/2020 e Resolução CFM nº 2.299/2021 \cite{cfm_digital:2021}.
    \end{itemize}

    \item \textbf{Dupla Checagem para Fármacos de Alta Vigilância}:
    \begin{itemize}[noitemsep]
        \item \textit{Declaração}: A liberação e a administração à beira do leito de medicamentos de alta vigilância (quimioterápicos, eletrólitos concentrados, insulina intravenosa, opioides e sedativos profundos) exigem conferência dupla independente.
        \item \textit{Ação Sistêmica}: O módulo de aprazamento de enfermagem exige a autenticação sequencial de dois profissionais de enfermagem distintos via leitura biométrica ou credenciais individuais no leito antes de liberar o registro de administração do fármaco.
        \item \textit{Fundamentação}: Metas Internacionais de Segurança do Paciente da Organização Mundial da Saúde (OMS / Joint Commission).
    \end{itemize}

    \item \textbf{Integridade Cadastral e Protocolo de Paciente Desconhecido}:
    \begin{itemize}[noitemsep]
        \item \textit{Declaração}: O cadastro regular de pacientes exige consistência obrigatória do número do CPF com a base de dados oficial da Receita Federal do Brasil.
        \item \textit{Ação Sistêmica}: O sistema rejeita cadastros com dígito verificador inválido. Para casos de emergência em que o paciente dê entrada inconsciente ou sem documentos, o sistema permite acionar o protocolo ``Paciente Desconhecido'', gerando identificador provisório que impede a interrupção do atendimento de urgência.
        \item \textit{Fundamentação}: Portaria do Ministério da Saúde sobre o Cartão Nacional de Saúde (CNS).
    \end{itemize}

    \item \textbf{Validação Prévia de Cobertura de Operadora de Saúde (TISS)}:
    \begin{itemize}[noitemsep]
        \item \textit{Declaração}: Procedimentos eletivos, exames de alta complexidade e internações cobertas por convênios de saúde suplementar requerem token formal de autorização emitido pela respectiva operadora.
        \item \textit{Ação Sistêmica}: O agendamento definitivo de salas cirúrgicas e reservas eletivas de leito permanece com status pendente até que a transação eletrônica de elegibilidade e autorização retorne aprovação da operadora via barramento TISS.
        \item \textit{Fundamentação}: Diretrizes operacionais e regulatórias da Agência Nacional de Saúde Suplementar (ANS) \cite{ans_tiss:2023}.
    \end{itemize}

    \item \textbf{Imutabilidade do Prontuário Encerrado e Notas de Retificação}:
    \begin{itemize}[noitemsep]
        \item \textit{Declaração}: Uma vez homologada a alta médica, a transferência externa definitiva ou o atestado de óbito, o prontuário eletrônico do episódio correspondente é trancado em caráter permanente.
        \item \textit{Ação Sistêmica}: A aplicação revoga privilégios de escrita ou edição sobre o texto das evoluções clínicas originais. Esclarecimentos posteriores ou correções de fatos clínicos pretéritos só podem ser inseridos sob a forma de notas aditivas assinadas digitalmente, preservando o histórico integral das versões.
        \item \textit{Fundamentação}: Resolução CFM nº 1.821/2007 \cite{cfm_prontuario:2007}.
    \end{itemize}
\end{enumerate}

% ==============================================================================
% 5.3 REQUISITOS NÃO FUNCIONAIS
% ==============================================================================
\subsection{Requisitos Não Funcionais (RNF)}

Os Requisitos Não Funcionais estabelecem os critérios de qualidade técnica, desempenho, confiabilidade e interoperabilidade do Vitalis v2.0. A Tabela \ref{tab:rnf_vitalis} sintetiza as métricas de engenharia estipuladas para a infraestrutura hospitalar.

\begin{table}[H]
\centering
\small
\caption{Matriz Consolidada de Requisitos Não Funcionais do Vitalis v2.0}
\label{tab:rnf_vitalis}
\begin{tabularx}{\textwidth}{l >{\raggedright\arraybackslash}p{3.5cm} >{\raggedright\arraybackslash}p{4.2cm} >{\raggedright\arraybackslash}X}
\toprule
\textbf{Código} & \textbf{Categoria de Qualidade} & \textbf{Métrica / Alvo Mensurável} & \textbf{Mecanismo Arquitetural de Garantia} \\
\midrule
\textbf{RNF-001} & Desempenho e Vazão do Banco de Dados & Leitura $< 200\,\text{ms}$; escrita transacional $< 400\,\text{ms}$ sob $2.000$ conexões simultâneas. & Índices B-Tree parciais, particionamento temporal de tabelas do PEP e réplicas de leitura desacopladas. \\
\addlinespace
\textbf{RNF-002} & Latência de Inferência de IA & Tempo de resposta no percentil 95 ($\text{P95}$) $< 1{,}5\,\text{s}$; timeout estrito configurado em $2{,}0\,\text{s}$. & Pipeline de RAG com indexação vetorial pré-computada em FAISS/HNSW e microsserviços dedicados. \\
\addlinespace
\textbf{RNF-003} & Segurança e Proteção Criptográfica & Trânsito com TLS 1.3 (ECDHE); repouso com algoritmo AES-256 (GCM). & Rotação periódica de certificados, chaves assimétricas em HSM e criptografia nativa no subsistema de armazenamento. \\
\addlinespace
\textbf{RNF-004} & Disponibilidade e Tolerância a Falhas & Disponibilidade mensal mínima de $99{,}95\%$ (indisponibilidade máxima de $21{,}9$ minutos/mês). & Topologia multi-zona com \textit{failover} automatizado de banco de dados e balanceamento de carga ativo-ativo. \\
\addlinespace
\textbf{RNF-005} & Interoperabilidade em Saúde & Aderência integral aos padrões HL7 v2.5 e HL7 FHIR Release 4 estruturados em JSON. & Camada de adaptação de barramento de integração em conformidade com as especificações da RNDS do Ministério da Saúde. \\
\addlinespace
\textbf{RNF-006} & Explicabilidade e Auditoria de Modelos & Rastreabilidade de $100\%$ das chamadas de inferência de IA com registro de contexto e escores. & Persistência em base append-only contendo modelo, versão, temperatura, tokens consumidos e referências clínicas do RAG. \\
\addlinespace
\textbf{RNF-007} & Acessibilidade e Ergonomia Clínica & Conformidade com diretrizes WCAG 2.1 nível AA; taxa de erro operacional $< 1\%$ em rounds clínicos. & Contraste tipográfico elevado ($\ge 4{,}5:1$), navegação guiada por teclado e componentes de toque para tablets $\ge 48 \times 48\,\text{dp}$. \\
\bottomrule
\end{tabularx}
\end{table}

Abaixo são detalhadas as especificações de cada requisito não funcional:

\begin{enumerate}[label=\textbf{RNF-\arabic*:}, leftmargin=*]
    \item \textbf{Desempenho e Vazão de Banco de Dados}:
    Em operações de pronto atendimento e enfermarias, os tempos de resposta para consultas cadastrais e abertura do PEP não podem ultrapassar $200\,\text{ms}$ em 99\% dos casos. Transações de escrita com garantia estrita de ACID (Atomicidade, Consistência, Isolamento e Durabilidade) devem ser concluídas em tempo inferior a $400\,\text{ms}$ sob carga nominal de $2.000$ requisições simultâneas. O banco de dados PostgreSQL adota conexões multiplexadas com PgBouncer, particionamento mensal das tabelas de auditoria e réplicas de leitura para relatórios e consultas do dashboard.

    \item \textbf{Latência de Inferência de IA}:
    As chamadas direcionadas aos modelos de linguagem clínica e mecanismos de RAG (\textit{Retrieval-Augmented Generation}) \cite{lewis_rag:2020} devem responder com latência no percentil 95 ($\text{P95}$) inferior a $1{,}5\,\text{s}$. Para assegurar que o profissional não enfrente congelamento de tela durante o ato médico, a camada de aplicação aplica timeout rígido de $2{,}0\,\text{s}$. Caso o modelo não finalize a inferência nesse período, o \textit{Circuit Breaker} encerra a chamada e aciona os algoritmos determinísticos de contingência.

    \item \textbf{Segurança e Criptografia em Trânsito e Repouso}:
    Toda comunicação interna e externa da arquitetura trafega obrigatoriamente através do protocolo TLS 1.3 com criptografia de ponta a ponta, suportando exclusivamente conjuntos de cifras que implementem sigilo de encaminhamento perfeito (\textit{Perfect Forward Secrecy}). Os volumes de armazenamento do banco de dados relacional, o armazenamento de laudos e imagens radiológicas e os snapshots de backup utilizam criptografia em repouso AES com chaves de 256 bits, gerenciadas por Módulo de Segurança de Hardware (HSM).

    \item \textbf{Disponibilidade e Tolerância a Falhas}:
    O sistema Vitalis v2.0 opera com nível de disponibilidade mínima garantida de $99{,}95\%$ ao mês, o que admite tempo máximo acumulado de indisponibilidade não planejada de $21{,}9$ minutos por mês. A infraestrutura baseia-se em cluster de orquestração de contêineres distribuído em pelo menos três Zonas de Disponibilidade (AZs) distintas, dotado de autoescalonamento horizontal e transição automática (\textit{failover}) de banco de dados com perda de dados nula ($\text{RPO} = 0$) para transações confirmadas.

    \item \textbf{Interoperabilidade e Padrões em Saúde}:
    Para viabilizar a comunicação com sistemas de laboratório (LIS), sistemas de arquivamento de imagens (PACS) e o barramento da Rede Nacional de Dados em Saúde (RNDS), a plataforma implementa suporte nativo aos padrões HL7 v2.5 e HL7 FHIR Release 4 \cite{hl7_fhir:2019}. Recursos como \texttt{Patient}, \texttt{Encounter}, \texttt{Condition}, \texttt{MedicationRequest} e \texttt{Observation} são expostos por endpoints RESTful estruturados em JSON.

    \item \textbf{Explicabilidade e Auditoria de Inferências de IA}:
    Em cumprimento aos princípios éticos da inteligência artificial médica, todas as recomendações emitidas pela IA clínica devem ser acompanhadas de metadados de rastreabilidade. Para cada resposta apresentada na interface, o sistema armazena e disponibiliza para auditoria: identificador e versão do modelo, parâmetro de temperatura, quantitativo de tokens consumidos, percentual de certeza estatística e trechos textuais dos artigos científicos e diretrizes farmacológicas recuperados pela busca vetorial no RAG.

    \item \textbf{Acessibilidade e Ergonomia Clínica}:
    A interface do usuário foi projetada para minimizar a carga cognitiva sob condições de pressão assistencial, alcançando conformidade com as Diretrizes de Acessibilidade para Conteúdo Web (WCAG 2.1) no nível AA. A tipografia adota taxa de contraste igual ou superior a $4{,}5:1$ para textos regulares, atalhos de teclado para navegação sem necessidade de mouse e componentes táteis para tablets com dimensões mínimas de $48 \times 48\,\text{dp}$, mantendo o índice de erros operacionais abaixo de $1\%$ em simulações clínicas padronizadas.
\end{enumerate}

% ==============================================================================
% 5.4 MATRIZ DE RASTREABILIDADE E CONFORMIDADE
% ==============================================================================
\subsection{Matriz de Rastreabilidade e Conformidade Regulatória}

A Tabela \ref{tab:rastreabilidade_vitalis} mapeia o alinhamento das regras de negócio e requisitos do Vitalis v2.0 com os marcos legais e as diretrizes normativas da saúde no Brasil.

\begin{table}[H]
\centering
\small
\caption{Matriz de Rastreabilidade entre Requisitos, Diretrizes e Normas Regulatórias}
\label{tab:rastreabilidade_vitalis}
\begin{tabularx}{\textwidth}{>{\raggedright\arraybackslash}p{2.6cm} >{\raggedright\arraybackslash}p{3.5cm} >{\raggedright\arraybackslash}p{3.5cm} >{\raggedright\arraybackslash}X}
\toprule
\textbf{Norma / Marco} & \textbf{Escopo Regulatório} & \textbf{Requisitos Vinculados} & \textbf{Regras de Negócio Associadas} \\
\midrule
CFM nº 1.821/2007 & Guarda, manuseio e descarte de prontuários eletrônicos. & RF-005, RF-009, RF-017 & RN-6 (Sigilo e Vínculo), RN-12 (Imutabilidade do Prontuário). \\
\addlinespace
CFM nº 2.299/2021 & Emissão de documentos médicos digitais e assinatura eletrônica. & RF-007, RF-009, RF-010 & RN-2 (Human-in-the-Loop), RN-8 (Validade ICP-Brasil). \\
\addlinespace
LGPD (Lei nº 13.709/18) & Proteção de dados pessoais sensíveis de saúde. & RF-001, RF-018, RF-019 & RN-3 (Desidentificação de Dados), RN-6 (Privilégio Mínimo). \\
\addlinespace
ANS (Padrão TISS) & Intercâmbio eletrônico de dados em saúde suplementar. & RF-012, RF-013 & RN-5 (Vedação Retroativa), RN-11 (Autorização TISS). \\
\addlinespace
Protocolo Manchester & Triagem clínica e estratificação de risco em urgência. & RF-002, RF-003, RF-021 & RN-1 (Prioridade Vermelha Absoluta). \\
\addlinespace
HL7 FHIR Release 4 & Interoperabilidade semântica e sintática de dados em saúde. & RF-001, RF-022, RNF-005 & RN-10 (Integridade Cadastral). \\
\bottomrule
\end{tabularx}
\end{table}

A consolidação estruturada dos requisitos funcionais, regras de negócio e requisitos não funcionais estabelece os critérios objetivos necessários para as etapas de projeto arquitetural, diagramação de processos e implementação técnica dos microsserviços do Vitalis v2.0.
